\documentclass[11pt, a4paper]{article}

% ── Encoding & fonts ──────────────────────────────────────────────────────────
\usepackage[utf8]{inputenc}
\usepackage[T1]{fontenc}
\usepackage{lmodern}

% ── Page geometry ─────────────────────────────────────────────────────────────
\usepackage[margin=1in]{geometry}

% ── Hyperlinks & auto-updating TOC ────────────────────────────────────────────
\usepackage[hidelinks]{hyperref}   % hidelinks = no ugly red boxes

% ── Math ──────────────────────────────────────────────────────────────────────
\usepackage{amsmath, amssymb, amsthm}

% ── Misc useful packages ──────────────────────────────────────────────────────
\usepackage{tabularx}              % equal-width wrapping columns
\usepackage{enumitem}              % tighter lists
\usepackage{parskip}               % space between paragraphs instead of indent

% ── Theorem-like environments (optional, delete if not needed) ─────────────────
\newtheorem{theorem}{Theorem}[section]
\newtheorem{lemma}[theorem]{Lemma}
\theoremstyle{definition}
\newtheorem{definition}[theorem]{Definition}
\newtheorem{example}[theorem]{Example}
\theoremstyle{remark}
\newtheorem*{remark}{Remark}

%tikz
\usepackage{tikz}
\usetikzlibrary{arrows.meta, positioning, fit, backgrounds}


% Code blocks
\usepackage{listings}
\usepackage{xcolor}

\lstset{
    basicstyle=\ttfamily\small,
    keywordstyle=\color{blue},
    commentstyle=\color{gray},
    stringstyle=\color{teal},
    numbers=left,
    numberstyle=\tiny\color{gray},
    numbersep=8pt,
    breaklines=true,
    frame=single,
    captionpos=b,
}

\lstnewenvironment{cppcode}[1][]{\lstset{language=C++, #1}}{}
\lstnewenvironment{pycode}[1][]{\lstset{language=Python, #1}}{}
\lstnewenvironment{ccode}[1][]{\lstset{language=C, #1}}{}
\lstnewenvironment{bashcode}[1][]{\lstset{language=bash, #1}}{}
\newcommand{\code}[1]{\lstinline|#1|}

% ── Chapters & appendices in any order ────────────────────────────────────────
% \chap{<n>}{Title}, \app{<n>}{Title} (1=A, 2=B, ...)
\newcommand{\chap}[2]{%
  \renewcommand{\thesection}{\arabic{section}}%
  \renewcommand{\theHsection}{chap.\arabic{section}}%
  \setcounter{section}{#1}\addtocounter{section}{-1}%
  \section{#2}}
\newcommand{\app}[2]{%
  \renewcommand{\thesection}{\Alph{section}}%
  \renewcommand{\theHsection}{app.\Alph{section}}%
  \setcounter{section}{#1}\addtocounter{section}{-1}%
  \section{#2}}


% ── Document metadata ─────────────────────────────────────────────────────────
\title{Computer Architecture: A Quantitative Approach}
\author{Hennessy \& Patterson}
\date{\today}

% ──────────────────────────────────────────────────────────────────────────────
\begin{document}

\maketitle
\tableofcontents          % auto-updates on every compile (needs 2 passes)
Order: Ch1, Appendix A, Appendix C, Appendix B, Ch2, Ch3, Ch5, Ch4, Ch7, Ch6
\newpage
\chap{1}{Fundamentals of Quantitative Design and Analysis}
\subsection{Introduction}
\begin{itemize}
    \item RISC - Reduced Instruction Set Computer was introduced in the early 1980s, with the aim of simplifying architectures to keep up with the rapid development the hardware.
    \item Utilized instruction level parallelism (initially with pipelining) and the use of caches.
    \item ARM became dominant in the late 90s as a RISC architecture, specifically for low-end power and cost sensitive applications.
\end{itemize}
\subsection{Classes of Computers}
\begin{itemize}
    \item There exist two main types of parallelism:
    \begin{enumerate}
        \item Data level parallelism (DLP) when there are many data items that may be parallelized.
        \item Task level parallelism (TLP) when larger tasks can operate independent of each other and may be parallelized.
    \end{enumerate}
    \item Hardware can then exploit these types of parallelism in four ways:
    \begin{enumerate}
        \item Instruction-level Parallelism (ILP) exploits DLP with the help of the compiler with ideas like pipelining and speculation.
        \item Vector Architectures and Graphic Processing Units (GPUs) exploit DLP by applying a single instruction to multiple pieces of data in parallel.
        \item Thread-Level Parallelism exploits either DLP or TLP in a specific hardware model that allows faster and efficient communication between threads
        \item Request-Level Parallelism exploits TLP as specified by the OS or programmer.
    \end{enumerate}
    \item Michael Flynn in the 60s all computers one of the following four categories sorted by parallelism in the instruction and data streams:
    \begin{enumerate}
        \item \textit{Single instruction, single data} (SISD) - A uniprocessor. Instructions are executed sequentially, but can still exploit ILP.
        \item \textit{Single instruction, multiple data} (SIMD) - There are multiple processors with independent data streams executing the same instruction in parallel. Each processor has its own data memory but they share the same instruction memory.
        \item \textit{Multiple instruction, single data} (MISD) - Doesnt really exist
        \item \textit{Multiple instruction, multiple data} (MIMD) - Targets TLP as each processor can fetch their own instructions and data. Generally more flexible and applicable than SIMD, but is inherently more expensive.
    \end{enumerate}
\end{itemize}
\subsection{Defining Computer Architecture}
\begin{itemize}
    \item Instruction set architecture (ISA): the set of instructions for the computer that are visible to the programmer. It serves as the boundary between software and hardware.
    \begin{enumerate}
        \item Class of ISA: Nearly all modern ISAs operate on registers or memory locations. 80x86 has 16 general-purpose regs and 16 floating-point regs, whereas MIPS has 32 of each.
        \begin{itemize}
            \item Register-memory ISA: Can access memory with a variety of instructions (eg 80x86)
            \item Load-store ISA: Can only access memory via load/store instructions (eg ARM and MIPS)
        \end{itemize}
        \item Memory addressing: Nearly all architectures use byte addressing to identify memory locations. Some architectures require an address to be aligned to a certain multiple (like ARM or MIPS) while others do not, but may have a performance cost.
        \item Addressing modes: 
        \begin{itemize}
            \item MIPS: Register, Immediate (constant only), Displacement
            \item 80x86: Same with three variants of displacement - no register (absolute), two registers (based indexed with displacement), and two registers (based with scaled index and displacement) 
            \item ARM: Same as MIPS plus PC-relative addressing, sum of two registers, and the sum of two registers where one register is multiplied by the size of the operand in bytes.
        \end{itemize}
        \item Types and sizes of operands: All 3 support operands sizes of 8-bit ascii, 16-bit unicode/half-word, 32-bit integer/word, 64-bit double word/long, 32-bit IEEE 754 floating point @ single precision, and 64-bit IEEE 754 floating point @ double precision. 80x86 also supports 80-bit extended floating point @ extended double precision.
        \item Operations: Categories include data transfer, arithmetic logical, control, and floating point.
        \item Control flow operations: Nearly all ISAs support conditional branching, unconditional jumps, procedure calls, and returns. Address is relative to program counter (PC). 
        \item Encoding an ISA: fixed length (eg ARM and MIPS, all instructions are 32 bits long) or variable length (eg 80x86, instructions are 1-18 bytes long). Variable length binaries tend to be smaller as instructions can be packed together more.
    \end{enumerate}
    \item Organization (or microarchitecture) refers to the high-level abstracted representation of a computer's design, such as the memory system, memory interconnect, and the design of the internal processor or CPU. 
    \item Two processors may implement the same ISA but have different organizations (eg AMD Opteron vs Intel Core i7).
    \item The hardware is the lower level specifics of a computer. 
    \item Two processors may implement the same ISA and organization but vary in physical hardware (eg Intel Core i7 vs Intel Xeon 7560).
    \item \textit{Architecture} is an overarching term that refers to the instruction set architecture, organization/microarchitecture, and hardware.
\end{itemize}
\subsection{Trends in Technology}
\begin{itemize}
    \item Well-designed ISAs are created with the intention to survive more rapid changes in technology. 
    \begin{enumerate}
        \item Integrated circuit logic technology: Transistor density (historically) has increased by 35\% each year, approximatley quadrupling over 4 years. Increases in die size range from 10\%-20\%, and the transistor count on a chip doubles every approx. 18-24 months (Moore's law)
        \item Semiconductor DRAM: Capacity per DRAM chip has increased by 25\%-40\% each year. This is slowing down however due to the difficulty of efficiently manufacturing small DRAM cellls.
        \item Semiconductor flash (electrically erasable programable read-only memory): Capacity per flash chip has increased ~50\%-60\% per year recently, and is 15-20 times cheaper than DRAM.
        \item Magnetic disk technology: Density has increased ~40\% per year since 2004. 15-25 times cheaper than flash and 300-500 times cheaper than DRAM, per bit.
        \item Network technology: Depends both on the performance of switches and the performance of the transmission system (hard to estimate)
    \end{enumerate}
    \item Bandwidth or Throughput: total amount of work done in a given time
    \item Latency or Response Time: Time between start and completion of a given task
    \item Generally, bandwidth outpaces latency improvements; generally bandwidth grows by at least the square of latency.
    \item Feature Size: minimum size of a transistor in either dimension. 
    \item Density of transistors increases quadratically wrt. a decrease in feature size
    \item Generally, transistor performance increases linearlly wrt. a decrease in feature size
    \item As feature size shrinks, however, wire performance degrades as they get shorter and the resistance and capacitance per unit length gets worse.
\end{itemize}
\subsection{Trends in Power and Energy in Integrated Circuits}
\begin{itemize}
    \item Three primary concerns from the perspective of a system designer:
    \begin{enumerate}
        \item What is the maximum power a processor requires?
        \item What is the sustained power consumption? Quantified via thermal design power (TDP), which determines the cooling requirement. Peak power is approx. 1.5 times TDP
        \item Energy efficiency
    \end{enumerate}
    \item Dynamic energy: energy required to switch a transistor in CMOS chips.
\end{itemize}
\begin{align*}
    \text{Energy}_{\text{dynamic}} \propto \text{Capacitive Load} \times \text{Voltage}^2
\end{align*}
Represents the energy pulse of the transition of 0$\rightarrow$1$\rightarrow$0 (or 1$\rightarrow$0$\rightarrow$1). Thus the energy of a single transition is then
\begin{align*}
    \text{Energy}_{\text{dynamic}} &\propto \frac{1}{2} \times \text{Capacitive Load} \times \text{Voltage}^2, \\
    \text{Power}_{\text{dynamic}} &\propto \frac{1}{2} \times \text{Capacitive Load} \times \text{Voltage}^2 \times \text{Frequency Switched}
\end{align*}
\begin{itemize}
    \item As technology evolves the increase of the number of transistors and frequency makes the decrease of load capacitance and voltage negligable, creating an overall increase of power consumption.
    \item Early microprocessors used $<1$ watt; modern processors consume on the order of 100s of watts.
    \item Power consumption has largely plateaued, as the issue of cooling the chips and preventing hot spots becomes increasingly more difficult.
    \item Various techniques exist to try and improve efficiency:
    \begin{enumerate}
        \item Do nothing: Turn off the clock on certain inactive modules
        \item Dynamic Voltage-Frequency Scaling (DVFS): When not operating at full capacity the clock rate will be lowered.
        \item Overclocking: Newer processors may allow certain cores to be able to operate at higher clock speeds for short periods of time, before having to cool off.
    \end{enumerate}
\end{itemize}
\subsection{Trends in Cost}
\begin{itemize}
    \item While the cost of manufacturing an IC has decreased dramatically, the overall process has not changed: a wafer gets tested and cut into dies. 
    \item Die yield: the fraction of good dies on a wafer relative to total dies. Empirically,
\end{itemize}
\begin{align*}
    \text{Die Yield} = \text{Wafer Yield} \times \frac{1}{(1+\text{Defects per Unit Area} \times \text{Die Area})^N}
\end{align*}
\begin{itemize}
    \item Wafer yield refers to the number of dies that are so bad they need not be tested; for simplicity this can be assumed to be \textasciitilde$100\%$.
    \item In 2010, Defects per unit area was reported to be 0.016-0.057 defects per square centimeter.
    \item $N$ is called the process-complexity factor, a measure of manufacturing difficulty. for 40 nm in 2010, this ranged from 11.5-15.5.
    \item To accomodate inevitable defects in manufacturing, many designers introduce redundant cells in the dies such that small defects still yield a functional die. This significantly increases the yield.
\end{itemize}
\subsection{Dependability}
\begin{itemize}
    \item Historically, ICs have been one of, if not the most reliable part of a computer. 
    \item While pins may be fragile the overall error rate inside the chip is very low.
    \item Service Level Agreement/Service Level Objective (SLA/SLO): A guarantee given by an infrastructure provider that their service would be dependable. 
    \item Two states: Service Accomplishment (delivered as specified) or Interruption (different from SLA). 
\end{itemize}
\subsection{Measuring, Reporting, and Summarizing Performance}
\subsection{Quantitative Principles of Computer Design}
\begin{itemize}
    \item Take advantage of parallelism: various techniques exist at different scales where tasks and/or data can be parallelized.
    \begin{itemize}
        \item System level: Consider a web service; the handling of mulitple requests can be spread out across multiple processors to improve throughput. 
        \item Single processor: Pipelining is when the processor will overlap one instruction while an earlier one is being executed to reduce the wait time between them. The key insight is that not every instruction depends on the previous one, so executing instructions partially or fully in parallel may be possible.
        \item Digital design: Modern ALUs use carry-lookahead, which use parallelism to speed the process of computing sums from linear to logarithmic to the number of bits.
    \end{itemize}
    \item Principle of Locality: Keep things that are used often close.
    \begin{itemize}
        \item As a rule of thumb, a program will spend 90\% of its runtime on just 10\% of the code.
        \item Based off of this we can predict what data and instructions a program will use based off of what the program has already done.
        \item Temporal locality: recently accessed items are more likely to be accessed again in the future.
        \item Spacial locality: items with similar addresses tend to be referenced close together in time.
    \end{itemize}
\end{itemize}
\subsubsection{Amdahl's Law}
Amdahl's Law defines the speedup that can be gained by a particular feature.
\begin{align*}
    \text{Speedup} = \frac{\text{Performance for entire task using enhancement when possible}}{\text{Performance for entire task without using the enhancement}}
\end{align*}
As well,
\begin{align*}
    \text{Execution Time}_\text{new} = \text{Execution Time}_\text{old}\times ((1-\text{Fraction}_\text{enhanced})+\frac{\text{Fraction}_\text{enhanced}}{\text{Speedup}_\text{enhanced}})
\end{align*}
Where $\text{Fraction}_\text{enhanced}$ is the fraction of the program that benefits from the enhancement, and $\text{Speedup}_\text{enhanced}$ represents how much faster said portion of the program is. Thus, the overall speedup ratio is then
\begin{align*}
    \text{Speedup}_\text{overall} = \frac{\text{Execution Time}_\text{old}}{\text{Execution Time}_\text{new}}=\frac{1}{(1-\text{Fraction}_\text{enhanced})+\frac{\text{Fraction}_\text{enhanced}}{\text{Speedup}_\text{enhanced}}}
\end{align*}
\newpage
\app{1}{Appendix A: Instruction Set Principles}
\subsection{Introduction}
\subsection{Classifying Instruction Set Architectures}
\begin{itemize}
    \item Main types of internal storage:
    \begin{enumerate}
        \item Stack: Operands are implicitly on top of the stack
        \item Accumulator: One operand is implicitly on top of the accumulator
        \item General-purpose register: Only explicit operands
    \end{enumerate}
    \item Consider the expression \texttt{C = A + B}:
\end{itemize}
\begin{center}
\begin{tabular}{c c c c}
\hline
Stack & Accumulator & Register-memory & Register (load-store) \\
\hline 
\hline 
\texttt{Push A}&\texttt{Load B}&\texttt{Load R1,A}&\texttt{Load R1,A}\\
\hline
\texttt{Push B}&\texttt{Add B}&\texttt{Add R1,B}&\texttt{Load R2,B}\\
\hline 
\texttt{Add}&\texttt{Store C}&\texttt{Store R1,C}&\texttt{Add R3,R1,R2}\\
\hline
\texttt{Pop C}&&&\texttt{Store R3,C}\\
\hline
\end{tabular}
\end{center}
\begin{itemize}
  \item Types of GPR architectures:
  \item Register-register (load/store): arithmetic operations can only operate on registers; memory is fetched through load/store operations. Takes 0-3 operands. Very fast and low CPI. (eg RISC)
  \item Register-memory: arithmetic operations can have one register and one memory operand. Takes 0-2 operands. Higher CPI, but code density is much higher than register-register.
  \item Memory-memory: arithmetic operations can fetch both operands from memory. Much slower and pipelining is nearly impossible as memory ops can stall. Essentially extinct in modern day architectures
\end{itemize}
\subsection{Memory Addressing}
\begin{itemize}
  \item Big endian vs little endian: Consider the number \texttt{0x12345678}. If we were to store this in address \texttt{0x00}:
\end{itemize}
\begin{center}
\begin{tabular} {c c c}
\hline 
Memory Address & Big Endian & Little Endian \\
\hline 
\hline 
0x00 & 12 (MSB) & 78 (LSB) \\
0x01 & 34 & 56 \\
0x02 & 56 & 34 \\ 
0x03 & 78 & 12 \\
\hline
\end{tabular}
\end{center}
\begin{itemize}
  \item Note that endianess does not affect the order of the bits themselves, but the bytes.
  \item Big endian is easier to read (follows left $\rightarrow$ right), but little endian is a little easier for the processor as the ALU processes bytes from LSB $\rightarrow$ MSB. Sign checking is also easier on little endian.
  \item CPUs will fetch memory in units known as words. If an object is in a misaligned memory address (an access to a object of size $s$ bytes with byte address $A$ where $A$ mod $s \neq 0$) will require multiple aligned memory accesses.
  \item Addressing mode: How architectures specify the address of an object they will access. The address specified by the mode is called the effective address.
\end{itemize}
\begin{center}
% \hsize factors must sum to 4 (the number of X columns); the stretchy
% \extracolsep soaks up rounding error that otherwise causes an underfull hbox
\begin{tabularx}{\textwidth}{@{\hspace{\tabcolsep}\extracolsep{0pt plus 1fil}}>{\centering\arraybackslash\hsize=0.8\hsize}X
                             >{\centering\arraybackslash\hsize=0.8\hsize}X
                             >{\centering\arraybackslash\hsize=1.6\hsize}X
                             >{\centering\arraybackslash\hsize=0.8\hsize}X}
\hline
Addressing Mode & Example Instruction & Meaning & When Used \\
\hline 
\hline 
Register & \texttt{Add R3,R4} & \texttt{Regs[R4]} $\leftarrow$ \texttt{Regs[R4] + Regs[R3]} & When a value is in a register \\ 
\hline
Immediate & \texttt{Add R4,\#3} & \texttt{Regs[R4]} $\leftarrow$ \texttt{Regs[R4] + 3} & When a value is a constant \\ 
\hline
Displacement & \texttt{Add R4,100(R1)} & \texttt{Regs[R4]} $\leftarrow$ \texttt{Regs[R4] + Mem[100 + Regs[R1]]} & Accessing local variables\\ 
\hline
Register indirect & \texttt{Add R4,(R1)} & \texttt{Regs[R4]} $\leftarrow$ \texttt{Regs[R4] + Mem[Regs[R1]]} & Access through a pointer\\ 
\hline
Indexed & \texttt{Add R3,(R1+R2)} & \texttt{Regs[R3]} $\leftarrow$ \texttt{Mem[Regs[R1] + Regs[R2]]} & Array indexing\\ 
\hline
Direct/absolute & \texttt{Add R1,(1001)} & \texttt{Regs[R1]} $\leftarrow$ \texttt{Regs[R1] + Mem[1001]} & Static data access\\ 
\hline
\end{tabularx}
\end{center}
\subsection{Types and Sizes of Operands}
\begin{itemize}
  \item Generally, the type of an operand (integer, single-precision floating point, char, etc) determines the size of the operand.
  \item Integers are almost always represented as twos-compliment binary numbers.
  \item Twos compliment: negative numbers are represented by inverting all bits and adding 1 (MSB becomes sign bit)
  \item Some architectures may also represent the operands as decimal values, since certain decimal fractions (like 0.1) do not have a binary representation.
\end{itemize}
\subsection{Operations in the Instruction Set}
\begin{center}
\begin{tabularx}{\textwidth}{l X}
\hline
\textbf{Operator Type} & \textbf{Examples} \\ 
\hline
Arithmetic and logical & Integer operations: add, subtract, and, or, multiply, divide, etc. \\
\hline 
Data transfer & Loads and stores (mov instructions on architectures with memory addressing) \\
\hline 
Control & Branch, jump, procedure call and return, traps \\
\hline 
System & syscall, virtual memory management instructions \\ 
\hline 
Floating-point & Floating-point operations: add, multiply, divide, compare \\ 
\hline 
Decimal & Decimal add, multiply, decimal->char conversion \\
\hline 
String & String move, cmp, search \\ 
\hline
Graphics & Pixel and vertex, compression/decompression operations \\
\hline
\end{tabularx}
\end{center}
\begin{itemize}
  \item The first 4 (Arithmetic, and logical, data transfer, control, and system) are the important ones 
  \item The top 10 most common instructions are responsible for 96\% of all instructions executed (on 80x86):
  \item 22\% Load, 20\% Conditional branch, 16\% Compare, 12\% Store  
\end{itemize}
\subsection{Instructions for Control Flow}
\begin{itemize}
  \item 4 types of control flow change:
  \begin{enumerate}
    \item Conditional branch
    \item Jump 
    \item Procedure call 
    \item Procedure return
  \end{enumerate}
  \item Majority is Conditional branch (\textasciitilde 80\%)
\end{itemize}
\subsubsection{Addressing modes for control flow instructions}
\begin{itemize}
  \item The destination address must always be specified.
  \item In the majority of cases, this is specified explicitly in the instruction
  \item Procedure return is the main exception, as the return address is not known at compile time (just pop call stack)
  \item Most common mode is via displacement relative to the program counter (PC). 
  \item Instructions that address this way are known as pc-relative instructions 
  \item Advantages are that these instructions take less bytes as the location is close to the current instruction.
  \item Position independence: PC-relative instructions also allow code to be run independent of the absolute location.
  \item Position independence helps with linking/dynamic linking.
  \item When PC-relative addressing is not available, the next best option is usually register indirect.
  \item This is due to the fact that PC-relative is not possible when the address is not known at compile time. 
  \item Register-indirect jumps also allow the following:
  \begin{enumerate}
    \item Switch-case functions
    \item Runtime dynamic dispatch
    \item Function pointers 
    \item Dynamically shared libraries
  \end{enumerate}
\end{itemize}
\subsubsection{Conditional branch options}
Below are ways in which conditional branches are evaluated:
\begin{enumerate}
  \item Condition code (CC): Branch operations will read specific flags set by the ALU (as a side effect) and branch accordingly. This has the advantage that it may be free as the arithmetic being done may have set the flag already. The disadvantage is that this depends on the state of the ALU flags, and can easily be overwritten restricting reordering of instructions. As well, the outcome is passed through the flags register to the branch, which is implicit (more dependencies for the compiler to consider) and there is usually only one flag register slot. Many modern architectures like x86 intel/amd, aarch64, and powerPC use CC.
  \item Condition register/limited comparison: A compare instruction will write a result into any GPR, which is then checked against a predicate. This doesn't rely on any flag states and is much simpler for the compiler to reason about. The disadvantage is that some conditions (eg a < b) introduce an extra compare instruction before the branching itself. This introduces a design tension, where either the branch's compare can be trivial (which introduces extra instructions) OR the branch's compare can be more powerful (but affect the critical path).
  \item Compare and branch: The compare operation is folded into the branch instruction itself, thus no extra instructions/flags are required, at the cost of a more expensive instruction.
\end{enumerate}
\end{document}
